\section*{5. Eingrenzung des Themas}
Der Fokus der Arbeit liegt darauf, zu analysieren, wie geeignet die Methode des Dashboards zur Modernisierung im Legacy-Kontext sein kann. Dafür wird untersucht, ob die Effizienz nach Einführung eines Dashboards, im Rahmen des Testpoolmonitorings, signifikant gesteigert werden kann.

Zur Einordnung der Qualität des erstellten Dashboards wird ein Usability-Test durchgeführt, dessen Ergebnis mithilfe einer Nutzerbefragung zusätzlich validiert wird. Die Ergebnisse beider Validierungsverfahren werden in den theoretischen Rahmen eingeordnet und dienen der Bewertung dieser Modernisierungsmaßnahme im Legacy-Kontext.

In dieser Untersuchung wird die Methode des Dashboards anhand einer Lösung im Monitoring-Kontext analysiert. Das bedeutet, dass weitere Methoden und Strategien zur Qualitätssteigerung und Prozessoptimierung sowie alternative Verfahren zur Ergebnisanzeige in diesem Kontext nicht weiter betrachtet werden.

Im Projektkontext ist das Dashboard-Tool auf die nach vorheriger Evaluation getroffenen Auswahl beschränkt. Nach Gewährleistung der benötigten Funktionalität, ist die Bewertung der dargestellten Inhalte, wie Vollständigkeit und Sinnhaftigkeit der Tests, nicht Teil der Arbeit. Weiterhin wird zur Untersuchung der Effizienz und Eignung des erstellten Dashboards, im Validierungsprozess ausschließlich das Entwicklungsteam einbezogen, welches bereits Zugriff auf das bestehende System besitzt und das Dashboard nutzen und warten wird.
